\section{从长乐宫到轮台：西汉怎样在自己的成功中变重}
\label{sec:han-late-western-depth}

长乐宫的门不会因为高祖去世而自动关上。前一九五年四月，刘邦死后，惠帝承接
皇位，吕后则承接了宫廷里更难量化的东西：谁能先接触新君，北军听谁的调动，
功臣的封国能否继续得到朝廷的答复。把这一阶段写成一个太后的私人故事，既会
误读传世叙事的道德化笔法，也会遮住正在运转的仓储、任命和军队。新皇帝不能
亲自逐一处理这些接口，太后与功臣也不能只靠血缘维持它们；他们必须让文书、
赏赐和守卫在次日继续走动。

\begin{event}{\eventanchor{WH-0195-GAOZU-DEATH}{高祖去世，惠帝即位}}{前一九五年四月}{汉 · 长安}{}
刘邦的死讯首先改变的，是长安哪些门可以通往皇帝。惠帝即位保住了刘氏继统，
却没有解除异姓功臣、宗室王国与禁军之间的紧张。吕后能够在太后的位置上介入
人事，正因为宫城里的保护、诏令和封爵需要有人把它们接到一起。对关东诸侯、
等待迁调的官员和供给京师的县邑而言，这不是宫闱边闻：他们要据此判断已有的
命令是否仍然有效，新的请求该向谁呈递。

《史记·吕太后本纪》《汉书·高祖纪下》和《惠帝纪》保存了继承的主干，却以
后来形成的汉代政治秩序安排人物。它们能证明谁即位、谁握有可见的职名，不能
让我们听见宫门内每一次谈判。较可靠的读法是把吕后看到的权力理解为一组必须
被维持的通道：她的处置会令这些通道更集中，也使后来反对她的人必须先争取
北军和宗室，才有可能改变局面。
\end{event}

惠帝死后的幼主，令这种集中更为明显。吕氏封王不是一张单纯的亲属名单；封国
意味着可供给的地方、可安插的官属和可能被调动的军政关系。反对者不必先宣称
一种抽象的政治原则，他们只需担心这些关系是否会永远排除刘氏宗室和旧功臣。
到吕后死后，周勃、陈平与刘氏诸王所进行的安排之所以能成功，正在于他们把
北军、宫门和继统的名义临时接在一起。诛吕并没有让军队从政治中消失；它只是
重新规定了哪一个网络能把军队称作“保卫汉室”。

文帝自代地入长安，提供的是一位可以使这套临时联盟暂时收束的宗室君主。关中
朝廷随即面对的却是更慢的恢复：战后家庭如何留下劳力，铜料和钱币怎样流通，
王国如何在不再直接挑战皇帝的条件下维持地方生活。前一七八年的减租诏与弛禁
铸钱，不能被合成一幅人人获益的图景。田租名目降低，仍须经县乡吏员辨认田地
和人户；铸钱的许可让铜料、工匠和商人得到新的空间，也使钱的轻重、成色和谁
能取得铜料成为新的争执。所谓“无为”的表面，实际是把一部分调节留给地方、
市场和家户，同时保留朝廷重新介入的能力。

\subsection{武帝的关中：新年号先抵达谁}

建元改元的文书会先在长安的官署、仪礼和任命中出现，然后才沿着郡国传递。
武帝在前一四一年即位时，继承的是削藩以后更易核验的王国、能够征调的关中
仓储，也继承了一个不能只靠节用应付的边防。赵绾、王臧的起落和五经博士的
设置，说明朝廷试图从礼制与经学中寻找可用的人才和语言；它们并不意味着天下
士人从此接受同一套知识。博士席位、师承和郡国举荐只是新的入口，入口的存在
会改变谁能被看见，却不能替地方学校、家学和不同经说规定同一速度。

前一三九年张骞西行，尤其能说明这类入口的边界。长安可以发给使者节杖，却不
能使他穿过匈奴控制的地带，更不能要求大月氏、乌孙或绿洲城邦按汉廷的利益
安排道路。张骞被留置和归来所带回的，是关于马匹、远方政权和路线的片段知识。
后来常被称为“丝绸之路”的联系，绝不能倒写为当时已经存在、等候汉军占有的
通道。使团每一次换马、投宿、转译和求见，都依赖路上不同的人愿不愿承担风险。

马邑的失败让这种限制变得尖锐。诱敌的交易与伏兵没有把单于带入城下，泄露的
计划反而使边市、居民和守备承担更直接的危险。朝廷随后选择更多主动出击，并
非一场失败必然推出的唯一策略；和亲、互市、守塞和远征本来就一直彼此竞争。
但当双方都不再相信旧有的交换可以控制袭扰，河西、朔方和北方草原就成为汉廷
必须决定是否持续付费的地带。

\begin{event}{\eventanchor{WH-0091-WITCHCRAFT}{巫蛊案冲进长安街巷}}{前九一年}{汉 · 长安}{}
巫蛊案使太子刘据与长安城内的军政通道在最坏的时刻相撞。告发、搜捕与诏令在
宫中生成，却必须通过城门、北军和官吏成为行动；当太子无法相信命令的来源，
而守军也无法判断谁代表皇帝时，武器和消息都落进了同一座城。传世材料保存了
江充、太子和武帝的戏剧性叙事，却不能让我们把每一段言辞或死亡当作逐时记录。

可以确知的是，这场危机没有只是改变一个家庭的命运。它暴露了内朝近臣、禁军
和告讼程序一旦失去相互核验，继承人也可能被迫在宫门之外寻找保护。太子败死
以后，皇帝仍能惩办涉案者、重新安置继统；但能使疑忌迅速变成内战的通道已经
留在国家内部。前八九年的轮台诏常被读成武帝晚年的悔悟，它更具体的意义是：
同一个朝廷同时承认了远征、屯田、转粮与内地劳力有无法再由胜利叙事遮盖的
成本。
\end{event}

\subsection{霍光的门槛与宣帝的尺度}

昭帝死于前七四年，长安不能等待一个自然成熟的继承人。霍光先迎昌邑王，又在
极短时间内废之，改立长期在民间成长的刘病已。废立奏议列出的过失，是为极端
行动制作合法性门槛的文字，不能逐项变作宫廷日记；它仍让人看到一件不安的
事实：遗诏、禁军和任命通道集中在辅政者手中时，皇帝的名号不足以保证皇帝
本人能够决定谁留在位上。

宣帝后来逐步收回人事主导，并不等于这套中介从此无害。赵充国在羌地提出屯田、
招抚和分化，选择的恰是比大规模推进更可持续、却会在边地留下更久驻军压力的
办法。盐铁会议则把同一问题放进关中官员、商人、农家和边费的争论里。两者
共同表明，西汉不能只问“强不强”：它必须问一条道路、一座塞城或一项专营由
谁付钱，谁承担劳役和价格的波动，谁又能在账册上把这种负担说成必要。

\begin{turningpoint}[西汉晚期的制度得失]
从吕后到霍光，中枢学会了在幼主、边费和继承突变时让官署不停摆；从武帝到
宣帝，它又学会以货币、专营、屯田和谈判延长自己的战略半径。所得是更强的
调度能力。代价是宫门、簿册和边军愈发成为少数人能够占住的接口。王氏后来
能够把辅政转成居摄，并非突然闯入一个完好的国家，而是接手了这些既有接口。
\end{turningpoint}
